% ==============================================================================
% T0/FFGFT Theory: Document 243
% Title: FFGFT in the Information Formalism — Mapping and Equivalence
% Author: Johann Pascher
% Date: May 2026
% Reference: Doc. 206 (Triangle-Matrix Reduction), Doc. 230 (Hilbert Space Translation),
%            Doc. 231 (Hilbert Space Extension), Doc. 133 (Fractal Correction)
% Occasion: IPI discussion with Onur Teker (UIFT) on information priority
% ==============================================================================

\section*{Preliminary Note}

This document examines whether and how the geometric description of FFGFT
can be mapped into an information formalism. The occasion is the discussion
with Onur Teker (UIFT) on whether geometry or information is the more
fundamental primitive.

The central thesis: \textbf{FFGFT is fully mappable into an information
formalism} --- just as Doc.~230 showed for the Hilbert space formalism.
However, this mapping is \emph{not an equivalence}: information is lost
under projection (winding numbers, $\theta_4$ phase, non-closure $\varepsilon$).
The mapping is a projection $\pi$, not an isomorphism.

\textbf{Methodological consequence:} That FFGFT can be mapped into an
information formalism does not prove that information is prior. It shows
that both descriptions are projections of the same structure --- with
different residuals.

% ==============================================================================
\section{Basic Structure of the Mapping}
% ==============================================================================

\subsection{The geometric state space in FFGFT}

In FFGFT the fundamental state space is the $T^4$ torus with winding
numbers $(n_\theta, n_\varphi) \in \mathbb{Z}^2$ and phase
$\theta_4 \in [0, 2\pi)$. A physical state $\Psi$ is fully
characterised by:

\begin{equation}
  \Psi \;=\; \bigl(n_\theta,\; n_\varphi,\; \theta_4,\; \varepsilon\bigr)
  \tag{243.1}
\end{equation}

where $\varepsilon \approx 0.001$ is the non-closure (Doc.~193).

\subsection{\texorpdfstring{Explicit definition of $\pi_I$}{Explicit definition of pi\_I}}

The projection $\pi_I$ is defined as:

\begin{equation}
  \pi_I:\; \Psi(n_\theta, n_\varphi, \theta_4, \varepsilon)
  \;\longmapsto\;
  \bigl(I(n_\theta, n_\varphi),\;\operatorname{sgn}(\theta_4),\;
        \langle\varepsilon\rangle_\text{therm}\bigr)
  \tag{243.2a}
\end{equation}

where $I = \log_2|\mathcal{W}|$ is the bit count of admissible winding
configurations, $\operatorname{sgn}(\theta_4)$ is the binary collapse
of the continuous phase, and $\langle\varepsilon\rangle_\text{therm}$
is the non-closure appearing only as a thermal expectation value.

\textbf{Important:} $|\mathcal{W}|$ must be finite for $I$ to be
well-defined. The T⁴ geometry provides this bound via the maximum
admissible winding number from the Planck-scale constraint
$n_\text{max} \sim 1/\xi$ (see Doc.~009). Without this geometric
bound, $I$ diverges --- demonstrating that the information count
\emph{presupposes} the geometric state space.

\subsection{Information-theoretic representation}

The winding state $(n_\theta, n_\varphi)$ defines a finite set of $N$
distinguishable configurations. The information content of a state is:

\begin{equation}
  I \;=\; \log_2 N \quad \text{bits}
  \tag{243.2}
\end{equation}

This is Boltzmann's formula in information-theoretic language:
$S = k_B \ln N$ corresponds to $I = \log_2 N$ bits with appropriate
scaling. The mapping:

\begin{align}
  (n_\theta, n_\varphi) &\;\longmapsto\; I(n_\theta, n_\varphi)
  = \log_2 |\mathcal{W}| \tag{243.3}
\end{align}

where $\mathcal{W}$ is the set of admissible winding configurations
defined by the $Z_3$ topology of $T^4$.

\subsection{What is lost in the mapping}

The information-theoretic projection $\pi_I$ discards:

\begin{itemize}
  \item The \textbf{continuous phase} $\theta_4$ --- collapsed to a
        binary distinction (phase present / absent)
  \item The \textbf{individual winding numbers} $(n_\theta, n_\varphi)$
        --- only their ratio $n_\varphi / n_\theta = 2s$ (spin)
        survives as a quantum number
  \item The \textbf{non-closure} $\varepsilon$ --- it appears in the
        information representation as ``noise'', not as a structural residual
\end{itemize}

This is the \emph{residual} of the projection $\pi_I$:
$R_I = \{\theta_4, (n_\theta, n_\varphi)_\text{full}, \varepsilon\}$.

% ==============================================================================

\subsection{Excursus: Extended Information Formalisms}

The limitations of $\pi_I$ are not absolute. Different information
formalisms have different residuals:

\begin{center}
\renewcommand{\arraystretch}{1.3}
{\footnotesize
\begin{tabular}{p{3.2cm}cccp{3.5cm}}
\toprule
\textbf{Formalism} & $\theta_4$ & $(n_\theta,n_\varphi)$ & $\varepsilon$ & \textbf{Residual} \\
\midrule
Classical-discrete (Shannon)
  & $\to$1 bit & $\to$spin & $\to$noise
  & $\{\theta_4, (n_\theta,n_\varphi)_\text{full}, \varepsilon_\text{geom}\}$ \\[3pt]
Continuous (diff.\ entropy)
  & Yes (density) & $\to$spin & Part.\ (Fisher)
  & $\{(n_\theta,n_\varphi)_\text{full}\}$ \\[3pt]
Quantum information (QIT)
  & Yes (phase) & $\to$spin & $\to$decoherence
  & $\{(n_\theta,n_\varphi)_\text{full}, \varepsilon_\text{geom}\}$ \\[3pt]
Topological (TDA)
  & $\to$bit & Yes (inv.) & $\to$noise
  & $\{\theta_4, \varepsilon_\text{geom}\}$ \\
\bottomrule
\end{tabular}
}
\end{center}

\textbf{Result:} No information formalism preserves all three quantities
$(\theta_4, (n_\theta,n_\varphi)_\text{full}, \varepsilon)$ simultaneously.
The T⁴ geometry is the only formalism that contains all three structurally.
The residual of $\pi_I$ is therefore \emph{formalism-relative} --- which
strengthens the core argument: which quantities are lost depends on the
chosen information formalism, not on the geometry.

\section{Concrete Translation Table}
% ==============================================================================

\begin{center}
\renewcommand{\arraystretch}{1.4}
{\small
\begin{tabular}{p{5.2cm}p{6.0cm}p{5.2cm}}
\toprule
\textbf{FFGFT (geometric)} & \textbf{Information formalism} & \textbf{Residual} \\
\midrule
Winding number $(n_\theta, n_\varphi)$
  & Discrete state, $\log_2 N$ bits
  & Individual values lost; only ratio preserved \\[4pt]
Phase $\theta_4 \in [0, 2\pi)$
  & 1 bit (phase present / absent)
  & Continuous information lost \\[4pt]
Non-closure $\varepsilon \approx 0.001$
  & Thermodynamic cost: $k_B T \ln 2$ per bit erasure
  & Geometric origin of $\varepsilon$ not visible \\[4pt]
$\tilde{T} \cdot m = 1$ (geometric, nat.\ units)
  & $\tau \cdot (\mu c^2) = \hbar$ (thermodynamic, SI);
    equiv.\ $\tau \cdot \mu = \hbar/c^2$ (SI: kg\,s)
    or $\tau \cdot \mu = 1$ (nat.\ units $\hbar=c=1$)
  & None --- this relation is invariant under $\pi_I$;
    $\tau\cdot\mu$ not dimensionless in SI \\[4pt]
$\xi = 4/30000$ (winding ratio)
  & $\xi = k_B T_\text{CMB} \ln 2 / (m_e c^2)$ (thermodynamic)
  & None --- $\xi$ survives as bridge parameter \\[4pt]
$\Delta\text{CHSH}(N) \approx 5 \times 10^{-4}$
  & Landauer cost per measurement: $\xi/(2\pi) \approx 2 \times 10^{-5}$
  & Integration scale lost (cumulative vs. elementary) \\[4pt]
Particle masses from windings
  & $m = r \cdot \xi^p \cdot v$ (formula)
  & Topological origin not visible \\[4pt]
\bottomrule
\end{tabular}
}
\end{center}

% ==============================================================================
\section{Mapping into a Neural Network}
% ==============================================================================

\subsection{Motivation}

A neural network is an information-processing system whose activation
patterns represent discrete states. If FFGFT states are representable
as activation patterns, this would be a concrete information-theoretic
realisation of FFGFT.

\subsection{Architecture}

A recurrent neural network (RNN) can implement the $\xi$-field equation

\begin{equation}
  \Psi(t) \;=\; G\bigl(\Psi(t - T_\text{rev}),\; \xi\bigr)
  \tag{243.4}
\end{equation}

as a recursion formula. The activation patterns of the hidden layer
correspond to winding states:

\begin{align}
  h_t &= \sigma\bigl(W \cdot h_{t-1} + b\bigr) \tag{243.5} \\
  \text{with}\quad W_{ij} &\sim (n_\theta, n_\varphi)_\text{coupling matrix}
  \tag{243.6}
\end{align}

The weight matrix $W$ encodes the topological couplings between winding
modes --- it is the information-theoretic counterpart of the T⁴ geometry.

\subsection{What the network cannot encode}

\begin{itemize}
  \item The \textbf{continuous phase} $\theta_4$: discrete activation
        functions collapse it
  \item The \textbf{exact non-closure} $\varepsilon$: it appears as
        training noise, not as a structural property
  \item The \textbf{geometric origin} of $\xi$: the network learns $\xi$
        as a hyperparameter, not as a winding ratio
\end{itemize}

The neural network is therefore a realisation of the projection $\pi_I$:
it correctly encodes the informationally accessible part of FFGFT,
but with the residuals noted above.

% ==============================================================================
\section{Relation to Doc.~230: Hilbert Space Translation}
% ==============================================================================

Doc.~230 showed that FFGFT is fully translatable into Hilbert space
quantum mechanics --- with the residual that winding numbers are lost.

The present document shows the same for the information formalism:
FFGFT is mappable --- with the residual that $\theta_4$, individual
winding numbers, and the geometric origin of $\varepsilon$ are lost.

\begin{center}
\renewcommand{\arraystretch}{1.3}
{\small
\begin{tabular}{p{4.5cm}p{2.5cm}p{1.5cm}p{4.5cm}}
\toprule
\textbf{Target} & \textbf{Projection} & \textbf{Equiv.?} & \textbf{Main residual} \\
\midrule
QM / Hilbert space & $\pi_\text{QM}$ & No & Winding numbers \\
SM (field theory) & $\pi_\text{SM}$ & No & $\theta_4$ phase \\
Information formalism & $\pi_I$ & No & $\theta_4$ + $\varepsilon$ origin \\
Matrix algebra & $\pi_\text{Mat}$ & No & Both winding numbers \\
\bottomrule
\end{tabular}
}
\end{center}

In all cases: FFGFT generates the target representation as a projection.
The target representation cannot fully reconstruct FFGFT.
\textbf{The mapping direction is asymmetric.}

% ==============================================================================
\section{Consequence for the Priority Debate}
% ==============================================================================

The mappability of FFGFT into the information formalism does \emph{not}
prove that information is the ontological prior. It shows:

\begin{enumerate}
  \item Both descriptions --- geometric and information-theoretic ---
        are internally consistent
  \item Both are projections of a common structure $\Psi$
  \item The geometric description has a smaller residual:
        it loses less when translated into other formalisms
  \item The parameter $\xi = 4/30000$ is the \emph{bridge parameter}
        that is identical in both descriptions --- as a winding ratio
        and as a thermodynamic constant $k_B T_\text{CMB} \ln 2 / (m_e c^2)$
\end{enumerate}

\textbf{Note on $\xi$-invariance:} The invariance of $\xi$ under $\pi_I$
is \emph{numerical}, not semantic: as the number $4/30000$ it is preserved,
but its meaning changes (winding ratio $\leftrightarrow$ thermodynamic
energy ratio). This dual meaning is precisely its bridge function.

\textbf{Working hypothesis:} The T⁴ geometry is not ``more true'' than
the information formalism, but it is \emph{derivationally richer}:
more physical parameters follow from it without additional input.
Whether this is an ontological statement or merely a methodological one
remains open --- deliberately, since physics operates at the level of
the measurable.


% ==============================================================================
\section{Is a Neural Network a Genuine Infor\-mation-Pro\-cessing System?}
% ==============================================================================

\subsection{The functional view}

At the functional level, a neural network is an information-processing
system in Shannon's sense: it receives inputs (distinguishable states),
transforms them through a sequence of operations, and produces outputs
(distinguishable states). Each layer performs a distinction --- between
activated and non-activated nodes --- and the composition of layers
produces a hierarchy of distinctions. This satisfies Shannon's
definition: $I = \log_2 N$ bits per resolved alternative.

A broader AI system --- combining neural networks with memory,
retrieval, symbolic reasoning, and language models --- extends this
further. It maintains structured state across queries, selects among
possible responses, and generates outputs that distinguish between
alternatives. In the functional sense, such a system processes
information in a well-defined way.

\subsection{The fundamental view: geometry as prior}

At the fundamental level, a neural network is a \emph{geometric
structure}: a directed graph with weighted edges (the weight matrix
$W$), nonlinear node functions (activation), and a fixed topology
(architecture). This geometric structure is defined entirely
\emph{before} any information is processed.

The architecture --- which nodes exist, how they are connected, what
the weight matrix looks like --- is a geometric object. The
information processing (the transformation of inputs to outputs) is
what happens when this geometric structure is \emph{instantiated} on
a physical substrate and evaluated.

This mirrors the FFGFT situation exactly:
\begin{itemize}
  \item T⁴ geometry defines which states are admissible
        (analogous to: the weight matrix defines which transformations
        are admissible)
  \item Physical particles are the instantiated states
        (analogous to: the network's output is the instantiated result)
  \item The information content follows from the geometry
        (analogous to: the number of distinguishable outputs follows
        from the network architecture)
\end{itemize}

\subsection{The Landauer argument revisited}

Onur Teker's argument is that Landauer's principle grounds
information-priority physically: erasing a bit dissipates $k_B T \ln 2$
as heat, independently of any observer. This is correct for a
\emph{physical} computation.

But consider where the Landauer cost arises in a neural network running
on a computer:
\begin{itemize}
  \item \textbf{Not} in the abstract weight matrix (a mathematical object)
  \item \textbf{Not} in the network architecture (a geometric structure)
  \item \textbf{Yes} in the transistor switching operations on the
        physical processor --- the substrate, not the network
\end{itemize}

The abstract neural network is geometrically defined and carries no
Landauer cost. The physical instantiation carries Landauer cost.
This is not a contradiction of Landauer's principle --- it is a
confirmation of the distinction between formal and physical information
that Doc.~190 R12 makes explicit for CHSH:

\begin{equation}
  \Delta\text{CHSH}_\text{formal} \;=\; 0
  \qquad\text{(no physical cost for abstract distinction)}
  \tag{243.8}
\end{equation}

\begin{equation}
  \Delta\text{CHSH}_\text{elem} \;=\; \frac{\xi}{2\pi}
  \;\approx\; 2\times10^{-5}
  \qquad\text{(Landauer cost per single physical measurement)}
  \tag{243.9}
\end{equation}

Note: the cumulative deviation $\Delta\text{CHSH}(N) \approx 5\times10^{-4}$
at $N=73$ runs (Doc.~022) is a \emph{different} observable from the
elementary cost $\xi/(2\pi)$ per measurement (Doc.~230). These are not
in contradiction; they describe the same physical effect at different
integration scales (Doc.~190 R12). The two values are consistent for
$N \approx 2\pi/\xi \approx 47{,}000$ --- the scale at which the
cumulative formula reaches the elementary value.

The same distinction applies to neural networks: the geometry is
formal; the Landauer cost arises when the geometry is physically
instantiated.

\subsection{Consequence for the priority debate}

A neural network is simultaneously:
\begin{itemize}
  \item An \textbf{information-processing system} at the functional
        level (Shannon)
  \item A \textbf{geometric structure} at the fundamental level
        (weight matrix, architecture, topology)
\end{itemize}

Neither description is more fundamental --- they are projections of
the same object at different levels of description. This is the
structure of $\tilde{T} \cdot m = 1$: the same physical fact viewed
from two sides. The geometry (weight matrix) is the condition; the
information processing (input-output transformation) is the function.
Neither generates the other; both are required for a complete
description.

\textbf{What AI systems show about the priority question:} Even the
most sophisticated AI --- combining language models, memory, retrieval,
and reasoning --- is at its foundation a geometric structure
(architecture, weights, connections) that is physically instantiated
on a substrate. The information processing is \emph{what the geometry
does when instantiated}. The geometry does not derive from the
information processing; the information processing derives from the
geometry being run.

This does not resolve the ontological question of which is prior. It
shows that the question is well-posed at every level of description ---
and that the geometric entry point is at least as natural as the
informational one, since every AI system begins with a geometric
design decision (architecture) before any information is processed.


% ==============================================================================
\section{Numerical Verification: RNN Weight Matrix}
% ==============================================================================

\subsection{The log-space network: geometry as a single weight}

The FFGFT mass formula $m = r \cdot \xi^p \cdot v$ is nonlinear in $p$,
so a standard linear network fails (errors of 50--2000\,\%). In
log-space it becomes exactly linear:

\begin{equation}
  \ln m \;=\; \underbrace{1}_{w_1}\cdot\ln r
           \;+\; \underbrace{\ln\xi}_{w_2}\cdot p
           \;+\; \underbrace{\ln v}_{b}
  \tag{243.11}
\end{equation}

The analytic weights require no training:
$w_1 = 1$, $w_2 = \ln\xi \approx -8.923$, $b = \ln v$.

\subsection{Results}

\begin{center}
\renewcommand{\arraystretch}{1.3}
{\small
\begin{tabular}{lcccc}
\toprule
\textbf{Particle} & $r$ & $p$
  & $m_\text{FFGFT}$ [eV] & $\Delta\,\%$ \\
\midrule
Electron & $4/3$ & $3/2$ & $5.054\times10^5$ & $1.09$ \\
Muon     & $16/5$ & $1$   & $1.050\times10^8$ & $0.57$ \\
Tau      & $25/9$ & $2/3$ & $1.785\times10^9$ & $0.46$ \\
\bottomrule
\end{tabular}
}
\end{center}

Experimental values: $m_e = 5.110\times10^5$\,eV,
$m_\mu = 1.057\times10^8$\,eV, $m_\tau = 1.777\times10^9$\,eV.

\subsection{What this shows}

The weight $w_2 = \ln\xi$ \emph{is} the T⁴ geometry --- encoded as
a single number. The network carries the \emph{what} ($\xi = 4/30000$);
it cannot carry the \emph{why} (the winding topology that enforces this
value). This confirms the projection $\pi_I$: functionally complete,
geometrically incomplete.

% ==============================================================================
\section{Quotient Structure: Why No Reverse Mapping}
% ==============================================================================

The mapping $\pi_I$ is not injective. Different T⁴ states can map to
the same information state:

\begin{equation}
  (n_\theta, n_\varphi) \;\sim\; (k\,n_\theta, k\,n_\varphi)
  \quad\text{for any } k \in \mathbb{Z}\setminus\{0\}
  \tag{243.10}
\end{equation}

Two winding configurations with the same \emph{ratio}
$n_\varphi/n_\theta$ produce the same spin $s = n_\varphi/(2n_\theta)$
and thus the same information state under $\pi_I$.
The information formalism is therefore a \textbf{quotient theory}
of FFGFT: it identifies configurations that are geometrically distinct.

This is not a deficiency --- it is the precise sense in which the
information formalism is a \emph{coarser} description. The residual
$R_I$ is not a bug but a feature: it marks exactly what the quotient
throws away.

\textbf{Consequence for priority:} An information formalism without
geometry cannot decide \emph{which} of the equivalent winding
configurations is physically realised. The T⁴ geometry provides the
\textbf{selection rule}: not all ratios $n_\varphi/n_\theta$ are
admissible --- only those satisfying the $Z_3$ topological constraint.
The information formalism says ``there are $N$ distinguishable states'';
the geometry says ``these specific $N$ are the admissible ones, and
here is why''.


% ==============================================================================
\section{Extension: Three AI Architectures as FFGFT Projections}
% ==============================================================================

Beyond RNNs, three fundamentally different architecture classes project
FFGFT differently and leave different residuals.

\subsection*{1. Transformer (GPT, BERT, Llama)}

Transformers use rotational position encodings (RoPE):
\begin{equation}
  \text{RoPE: } \mathbf{q}_m = e^{im\theta} \cdot \mathbf{q},
  \quad \theta \in [0, 2\pi)
  \tag{243.12}
\end{equation}
The phase $\theta_4$ can be encoded as relative position --- but as a fixed
encoding, not a dynamic variable. Residual: $\varepsilon$ is lost; $\xi$
is learned as attention temperature, not as winding ratio.

\subsection*{2. Modern Hopfield Network (Demircigil et al.\ 2017)}

The modern Hopfield network has an exponential energy function with
inverse temperature parameter $\beta = 1/T$:
\begin{equation}
  E = -\frac{1}{\beta}\log\!\left(\sum_\mu e^{\beta\cdot\text{sim}(\xi,\xi_\mu)}\right)
  \tag{243.13}
\end{equation}
For finite $\beta$, $\varepsilon$ appears as a structural parameter, not
training noise. \textbf{This is the best information-theoretic approximation
of $\varepsilon$ so far.} Residual: $\theta_4$ and individual winding numbers
are still lost.

\subsection*{3. Graph Neural Networks (GNNs)}

Winding numbers $(n_\theta,n_\varphi)$ can be encoded as weighted nodes.
Residual: $\theta_4$ and $\varepsilon$ are lost.

\subsection*{Summary of three architectures}

\begin{center}
\renewcommand{\arraystretch}{1.3}
{\small
\begin{tabular}{p{3.5cm}p{3cm}p{3cm}p{3cm}}
\toprule
\textbf{Architecture} & \textbf{$(n_\theta,n_\varphi)$?} & \textbf{$\theta_4$?} & \textbf{$\varepsilon$?} \\
\midrule
RNN (real) & Partial & No & No \\
Transformer (RoPE) & Indirect & Partial (fixed) & No \\
Hopfield (modern) & No & No & \textbf{Yes} (as $\beta$) \\
GNN & Yes (nodes) & No & No \\
\midrule
\multicolumn{4}{l}{\textbf{No} architecture encodes all three simultaneously.} \\
\bottomrule
\end{tabular}
}
\end{center}

% ==============================================================================
\section{The Inverse Problem: Can AI Reconstruct FFGFT?}
% ==============================================================================

\textbf{No-free-lunch for topological invariants:} From finitely many bits,
the complete $T^4$ geometry cannot be uniquely reconstructed:
\begin{equation}
  \nexists\;\text{AI},\;\text{reconstructing }
  (n_\theta,n_\varphi,\theta_4,\varepsilon)
  \text{ uniquely from }
  \{I,\Theta(\theta_4),\langle\varepsilon\rangle\}
  \tag{243.14}
\end{equation}

Current research (representation learning, persistent homology, contrastive
learning) can approximate individual aspects but not the specific number
$\xi = 4/30000$ --- it must be set as a hyperparameter and does not follow
from information alone.

% ==============================================================================
\section{Testable Prediction: AI Hardware as $T^4$}
% ==============================================================================

Every physical AI system runs on a processor with:
\begin{itemize}
  \item A clock (cyclic time: $\theta_4$ as clock phase)
  \item Finite registers (discrete states: $\log_2 N$ bits)
  \item Thermal noise ($\varepsilon$ as jitter)
\end{itemize}

\textbf{FFGFT prediction:} Torus-networked chips show lower Landauer costs
than 2D-grid chips:

\begin{center}
\renewcommand{\arraystretch}{1.3}
{\small
\begin{tabular}{p{3.5cm}p{1.5cm}p{3cm}p{4cm}}
\toprule
\textbf{Chip topology} & $\tau$ & \textbf{Exp.\ cost} & \textbf{FFGFT prediction} \\
\midrule
2D-Grid (standard) & 0 & $k_BT\ln 2$ & higher than $\xi$ \\
2D-Torus (TPU) & 0.5 & $\approx 0.67\cdot k_BT\ln 2$ & between $\xi$ and standard \\
3D-Torus (Cerebras) & 0.8 & $\approx 0.33\cdot k_BT\ln 2$ & closer to $\xi$ \\
4D-Torus (optical) & 1.0 & $\xi\cdot m_ec^2/\ln 2 = k_BT_\text{CMB}$ & exactly $\xi$ \\
\bottomrule
\end{tabular}
}
\end{center}

\textbf{Experimental decision of the priority debate:}
\begin{itemize}
  \item If costs depend on chip topology: geometry (torus) is fundamental,
        information follows from it.
  \item If all topologies show equal costs: information is fundamental.
\end{itemize}

FFGFT's prediction is clear: \textbf{Torus topologies win.}

% ==============================================================================

% ==============================================================================
\section*{Notation: Chip Topology Parameters}
% ==============================================================================

\begin{center}
\renewcommand{\arraystretch}{1.3}
{\small
\begin{tabular}{p{2cm}p{10cm}}
\toprule
\textbf{Symbol} & \textbf{Meaning} \\
\midrule
$\tau$ & Torus similarity of chip topology ($\tau=0$: 2D grid, $\tau=1$: ideal $T^4$) \\
$\alpha$ & Material-dependent coupling factor (photon-phonon vs.\ Coulomb) \\
OCS & Optical Circuit Switch (MEMS-based) \\
ICI & Inter-Chip Interconnect (Google proprietary high-speed protocol) \\
WSE & Wafer-Scale Engine (Cerebras) \\
$T^4$ & 4-torus: four compact loops (FFGFT full geometry) \\
$T^3$ & 3-torus: three compact loops (TPUv4 physical topology) \\
RoPE & Rotational Position Encoding (Transformer position encoding) \\
$\beta$ & Inverse temperature parameter in modern Hopfield network ($\beta = 1/T$) \\
$\pi_\text{elec}$ & Projection onto electrical mesh topology (additional residual) \\
\bottomrule
\end{tabular}
}
\end{center}


% ==============================================================================

\subsection{Clarification: Photonic vs.\ Electrical Chip Topologies}

\textbf{Not all torus chips are photonic.} The following overview
distinguishes between logical and physical topology realisation,
based exclusively on verified sources.

\begin{center}
\renewcommand{\arraystretch}{1.4}
{\footnotesize
\begin{tabular}{p{4.0cm}p{3.0cm}p{2.0cm}p{6.5cm}}
\toprule
\textbf{System} & \textbf{Topology} & \textbf{Phot.?} & \textbf{Source / Note} \\
\midrule
Google TPUv4 Pod
  & 3D Torus ($4{\times}4{\times}4$ cube)
  & \textbf{Yes} (OCS)
  & Jouppi et al.\ 2023, arXiv:2304.01433 \newline
    Palomar OCS: MEMS mirrors, $136{\times}136$ ports \newline
    Wrap-around via optical fibre, reconfigurable in 10\,s \\[4pt]
Google TPUv5p
  & 3D Torus ($16{\times}20{\times}28$)
  & \textbf{Yes} (OCS)
  & 13,824 optical ports, 48 OCS units \newline
    (based on Google documentation) \\[4pt]
Morphlux (Cornell/Lightmatter)
  & Reconfigurable ($T^3$ emulable)
  & \textbf{Yes} (silicon photonics)
  & Kumar et al.\ 2025, arXiv:2508.03674 \newline
    Programmable photonic chip-to-chip interposer \newline
    1.72$\times$ training throughput improvement \\[4pt]
Cerebras WSE-3
  & 2D Mesh (not a torus)
  & No (electrical)
  & Cerebras SDK documentation; sdk.cerebras.net \newline
    900,000 cores, 2D rectangular mesh on one wafer \\[4pt]
Standard GPU cluster
  & 2D Grid / Fat-Tree
  & No (electrical)
  & NVLink / InfiniBand between chips \\
\bottomrule
\end{tabular}
}
\end{center}

\textbf{Correction:}
Cerebras WSE-3 is not a torus but a \textbf{2D mesh} ---
topologically related but without the periodic boundary conditions
that define a torus. The FFGFT prediction (lower Landauer costs
for torus topologies) therefore applies to TPUv4/v5p,
not to the Cerebras WSE-3.

\subsubsection*{Why photonic torus chips are particularly relevant for FFGFT}

In the TPUv4, the optical circuit switch (OCS) realises the torus geometry
\emph{physically}: a light signal leaving the torus at one edge re-enters
at the opposite edge --- exactly like the $T^4$ boundary condition.
The light ``lives'' on the torus.

In electrical mesh topologies (Cerebras WSE-3), the connection is defined
logically by routing tables, not by the physical propagation of signals.
This corresponds to a projection $\pi_\text{elec}$ with additional
residuals compared to the ideal torus.

\textbf{Consequence:} The proposed experiment (table of $\tau$
vs.\ Landauer costs in Doc.~243 Section~12) should treat photonic
torus chips (TPUv4/v5p) and electrical mesh chips (Cerebras WSE-3)
separately. The limit $\tau \to 1$ is only expected for photonic 4D tori.


% ==============================================================================
\section{Open Points and Next Steps}
% ==============================================================================

\subsection*{Answered by Calculation}

\begin{itemize}
  \item \textbf{Numerical verification (answered):}
        The log-space network with a single weight $w_2 = \ln\xi \approx -8.923$
        reproduces all three lepton masses without training:
        electron $\Delta=1.09\,\%$, muon $\Delta=0.57\,\%$, tau $\Delta=0.46\,\%$.
        The network is $\pi_I$ in pure form: functionally complete, geometrically incomplete.

  \item \textbf{Constraint $T^+ + T^- = 0$ (partially answered):}
        An antisymmetric weight matrix $W = -W^\top$ can encode this condition
        structurally. For state $h = (h^+, h^-)$ with $h^+ + h^- = 0$,
        the condition is preserved as a conservation law when $W$ is antisymmetric
        and $h$ is antisymmetrically initialised. However, this requires a special
        architecture (antisymmetric weights) and is not present in a standard RNN.
\end{itemize}

\subsection*{Still Open}

\begin{itemize}
  \item \textbf{Formal derivation} of the network architecture from the
        T⁴ projection geometry (analogous to Doc.~189 for spinors)

  \item \textbf{Coefficient $275/4$} in the $T_\text{CMB}$ correction:
        $T_\text{CMB} = \tfrac{16}{9}\xi(1 - \tfrac{275}{4}\xi) = 2.725486\,\text{K}$
        ($\Delta = 0.0002\,\%$). The tetrahedral number $68 = 72-4$ (Doc.~003) gives
        $\Delta = 0.01\,\%$, the number $69$ gives $\Delta = 0.003\,\%$.
        The exact value $68.75 = 275/4$ has no geometric derivation yet.

  \item \textbf{Connection to UIFT}: If $\xi$ is measured thermodynamically
        (Vopson-Teker), this confirms the bridge --- but not the priority direction.
        Whether $\xi$ is primarily geometric or thermodynamic remains undecided
        by this test.
\end{itemize}

\subsection*{Literature (Chip Topologies)}

\begin{itemize}\small
  \item \textbf{Google TPUv4:} Jouppi, N.\,P.\ et al.\ (2023).
        TPU v4: An Optically Reconfigurable Supercomputer for Machine
        Learning with Hardware Support for Embeddings.
        \textit{Proc.\ ISCA 2023}. arXiv:2304.01433.
        Quote: \textit{``Optical circuit switches (OCSes) dynamically
        reconfigure its interconnect topology [\ldots]; users can pick
        a twisted 3D torus topology if desired.''}

  \item \textbf{Morphlux:} Kumar, A.\,V., Ding, E., Devraj, A.,
        Bunandar, D., \& Singh, R.\ (2025).
        Morphlux: Transforming Torus Fabrics for Efficient
        Multi-tenant ML. arXiv:2508.03674. Accepted at ASPLOS 2026.

  \item \textbf{Cerebras WSE-3:} Cerebras Systems.
        \textit{Cerebras SDK Documentation --- A Conceptual View}.
        sdk.cerebras.net. Quote: \textit{``The PEs are interconnected
        by communication links into a two-dimensional rectangular
        mesh on one single silicon wafer.''}
\end{itemize}


% ==============================================================================
\section*{Internal Document References (FFGFT Series)}
% ==============================================================================

The following internal documents are cited in the text. They are available
in the GitHub repository:
\texttt{github.com/jpascher/T0-Time-Mass-Duality}

\begin{center}
\renewcommand{\arraystretch}{1.3}
{\small
\begin{tabular}{p{1.2cm}p{6cm}p{6cm}}
\toprule
\textbf{Doc.} & \textbf{Title} & \textbf{Content (relevant for Doc.~243)} \\
\midrule
009 & Origin of $\xi$ & $K=245$, $\kappa=7$, e-p-$\mu$ triangle \\
133 & Fractal Correction Derivation & $K_\text{frak} = 1 - 100\xi$, RG run \\
189 & Spinor Winding Structure & $(n_\theta, n_\varphi)$ mapping to spin \\
190 & Corrections Register & R10--R12: $K=245$, $T_\text{CMB}$, CHSH \\
193 & Non-closure $\varepsilon$ & $\varepsilon \approx 0.001$ as arrow of time \\
206 & Triangle-Matrix Reduction & Geometry generates matrix, not vice versa \\
230 & Hilbert Space Translation & Bijective FFGFT $\leftrightarrow$ QM mapping \\
231 & Hilbert Space Extension & Bell states, CHSH from FFGFT \\
\bottomrule
\end{tabular}
}
\end{center}

\begin{keyresult}
  \textbf{Core statement of Doc.~243:}

  FFGFT is mappable into the information formalism (as projection
  $\pi_I$), analogous to the mapping into Hilbert space (Doc.~230).
  In the mapping, the $\theta_4$ phase, individual winding numbers,
  and the geometric origin of the non-closure are lost. The parameter
  $\xi$ is invariant under $\pi_I$ and serves as bridge between
  geometric and thermodynamic description. Mappability does not prove
  ontological priority of information --- it shows that both descriptions
  are projections of the same structure.
\end{keyresult}
